\section{Gauge fields: \texorpdfstring{Ba\l aban}{Balaban}'s gauge and coercivity}\label{sec:gauge-coercivity}

\subsection{The quadratic action and its gauge invariance}

The Wilson action of a lattice gauge field $U=e^{\ii\eta A}$ near $U=1$ is $\frac12\|dA\|^2$ plus terms of order three and higher in $A$ (\cite[(1.3)]{BalPropI}, \cite[(26)]{BalVar}). The quadratic part is invariant under gauge transformations
\[
A\mapsto A+d\lambda,
\]
with $\lambda$ an arbitrary scalar field. By Lemma~\ref{lem:avg}(d), $\bQ(A+d\lambda)=\bQ A+\dl Q_0\lambda$. Hence the constraint $\bQ A=B$ is preserved exactly when $\dl Q_0\lambda=0$, that is, when $Q_0\lambda$ is constant. Constant $\lambda$ act trivially, so the residual gauge group may be taken to be $\cN(Q_0)$. Consequently the variational problem
\[
\text{minimize }\|dA\|^2\quad\text{subject to}\quad\bQ A=B
\]
has, at best, a minimizer that is unique up to $A\mapsto A+d\lambda$ with $\lambda\in\cN(Q_0)$. A gauge condition is needed to select one representative.

\subsection{\texorpdfstring{Ba\l aban}{Balaban}'s gauge}

Ba\l aban uses the condition
\begin{equation}\label{eq:landau}
Rd^*A=0,
\end{equation}
with $R$ the orthogonal projection onto $\cW=-\Lap\cN(Q_0)$ (Proposition~\ref{prop:R}). He calls it the Landau gauge \cite[Sect.~D]{BalPropI}, \cite[(21)]{BalVar}.

\begin{lemma}[Gauge fixing]\label{lem:gauge-fix}
For every gauge field $A$ there is exactly one $\lambda\in\cN(Q_0)$ such that $A+d\lambda$ satisfies \eqref{eq:landau}. Moreover $\|d\lambda\|\le\sqrt{d/2}\,\|Rd^*A\|$, and $d(A+d\lambda)=dA$, $\bQ(A+d\lambda)=\bQ A$.
\end{lemma}

\begin{proof}
For $\lambda\in\cN(Q_0)$ we have $d^*d\lambda=-\Lap\lambda\in\cW$, hence $Rd^*d\lambda=-\Lap\lambda$. The condition $Rd^*(A+d\lambda)=0$ becomes $-\Lap\lambda=-Rd^*A$. The right-hand side lies in $\cW$, and $-\Lap$ is a bijection from $\cN(Q_0)$ onto $\cW$. The bound is Lemma~\ref{lem:kerQ0}.
\end{proof}

By Proposition~\ref{prop:R}(c), condition \eqref{eq:landau} says that $d^*A$ is a constant plus $(-\Lap)^{-1}Q_0^*\omega$ for some unit-lattice $\omega$. So $d^*A$ is the potential of a block-constant charge density. It is a $C^{1,\alpha}$ function on the unit scale, not a rough one. By contrast, the gauge that minimizes $\|A\|$ on the orbit would require $d^*A\in\ran Q_0^*$, a field with jumps at block boundaries. This regularity is one reason for Ba\l aban's choice.

\begin{remark}[Origin of the operator $\bL_a$]
Ba\l aban arrives at \eqref{eq:landau} by a Faddeev--Popov argument \cite[(1.21)--(1.39)]{BalPropI}. It replaces the axial gauge fixing by the density $\exp(-\frac12\|Rd^*A\|^2)$ and then takes the Landau limit. Adding the gauge-fixing term and the averaging term to the action gives the quadratic form
\[
\ip A{\bL_aA}=\|dA\|^2+\|Rd^*A\|^2+a\|\bQ A\|^2,\qquad
\bL_a=d^*d+dRd^*+a\bQ^*\bQ=-\Lap-dP_\cV d^*+a\bQ^*\bQ ,
\]
where we used Lemma~\ref{lem:hodge} and $R=I-P_\cV$. Up to the term $dP_\cV d^*$, which has exponential decay, this is a Feynman-type gauge: the principal part is the componentwise Laplacian. On the constraint set $\{\bQ A=B,\ Rd^*A=0\}$ the form reduces to $\|dA\|^2+a\|B\|^2$.
\end{remark}

\subsection{The averaging nearly separates exact and co-closed fields}

The heart of the coercivity proof is the following identity. On the unit lattice, gradients $\dl g$ and co-closed fields are orthogonal. We show that $\bQ$ maps fine co-closed fields to unit fields that are co-closed up to a divergence controlled by derivatives.

\begin{lemma}\label{lem:Q-divergence}
There is $C_W$, depending only on $d$, with the following property. For every gauge field $A$ on $\Tf$,
\[
\dl^*\bQ A=\iota\,b\,\bar c\,d^*A+\dl^*\mathsf WA,\qquad \bar c=c_1c_2\cdots c_d,
\]
where $\mathsf WA$ is a gauge field on $\T$ with $\|\mathsf WA\|\le C_W\|\partial A\|$.
\end{lemma}

\begin{proof}
By Lemma~\ref{lem:1d}(a), $D_\mu^*=a_\mu^*\partial_\mu^*$, and all operators involved commute. Hence
\[
\dl^*\bQ A=\sum_\mu D_\mu^*\iota\,a_\mu bA_\mu=\iota\sum_\mu b\,c_\mu\,\partial_\mu^*A_\mu .
\]
For fixed $\mu$, the telescoping identity $1-x_1\cdots x_m=\sum_ix_1\cdots x_{i-1}(1-x_i)$ for commuting $x_i$ gives
\[
c_\mu-\bar c=c_\mu\Big(1-\prod_{\nu\ne\mu}c_\nu\Big)=\sum_{\nu\ne\mu}c_\mu\Pi_{\mu\nu}(1-c_\nu),\qquad\Pi_{\mu\nu}=\prod_{\nu'<\nu,\ \nu'\ne\mu}c_{\nu'} .
\]
By Lemma~\ref{lem:1d}(d), $1-c_\nu=e_\nu\partial_\nu^*\partial_\nu$. By Lemma~\ref{lem:1d}(a),(b), $a_\nu\partial_\nu^*=\tau_\nu a_\nu^*\partial_\nu^*=\tau_\nu D_\nu^*$, so that $b\,\partial^*_\nu=\big(\prod_{\nu'\ne\nu}a_{\nu'}\big)\tau_\nu D^*_\nu$. Since $D^*_\nu$ commutes with $\iota$,
\[
\iota\,b\,(c_\mu-\bar c)\,\partial^*_\mu A_\mu=\sum_{\nu\ne\mu}D_\nu^*\,\iota\Big(\prod_{\nu'\ne\nu}a_{\nu'}\Big)\tau_\nu\,c_\mu\Pi_{\mu\nu}\,(e_\nu\partial_\nu)\,\partial^*_\mu A_\mu .
\]
Summing over $\mu$ and collecting the coefficient of $D^*_\nu$ defines
\[
(\mathsf WA)_\nu=\sum_{\mu\ne\nu}\iota\Big(\prod_{\nu'\ne\nu}a_{\nu'}\Big)\tau_\nu\,c_\mu\Pi_{\mu\nu}\,(e_\nu\partial_\nu)\,\partial^*_\mu A_\mu,
\]
and the identity follows. Each operator acting on $\partial^*_\mu A_\mu$ is $\iota$ composed with a product convolution. In direction $\nu$ it is $\tau_\nu e_\nu\partial_\nu$, whose kernel is bounded by $1$ and supported in an interval of length at most $3$ (Lemma~\ref{lem:1d}(d)). In every other direction it is $a_{\nu'}$ times at most one triangle average $c_{\nu'}$, whose kernel is an average of $c$ and therefore bounded by $1$, with support of length at most $3$. By Lemma~\ref{lem:restriction} the norm is at most $4^d\|\partial^*_\mu A_\mu\|=4^d\|\partial_\mu A_\mu\|$. Hence, by Cauchy--Schwarz over $\mu\ne\nu$, $\|\mathsf WA\|\le4^d d\,\|\partial A\|$.
\end{proof}

\begin{lemma}[Separation]\label{lem:separation}
There is $C_\star$, depending only on $d$, such that for every co-closed gauge field $B'$ on $\Tf$ and every scalar field $g$ on $\T$,
\[
\|B'\|^2+\|\dl g\|^2\le C_\star\big(\|dB'\|^2+\|\bQ B'+\dl g\|^2\big).
\]
\end{lemma}

\begin{proof}
Expand the square:
\[
\|\bQ B'+\dl g\|^2=\|\bQ B'\|^2+\|\dl g\|^2+2\ip{\dl^*\bQ B'}g .
\]
By Lemma~\ref{lem:Q-divergence} and $d^*B'=0$, $\ip{\dl^*\bQ B'}g=\ip{\mathsf WB'}{\dl g}$. This is bounded in absolute value by $C_W\|\partial B'\|\|\dl g\|\le\frac14\|\dl g\|^2+C_W^2\|\partial B'\|^2$. By Lemma~\ref{lem:hodge}, $\|\partial B'\|=\|dB'\|$. Hence
\[
\|\bQ B'+\dl g\|^2\ge\|\bQ B'\|^2+\tfrac12\|\dl g\|^2-2C_W^2\|dB'\|^2 .
\]
On the other hand, Corollary~\ref{cor:poincare-all} gives $\|B'\|^2\le C_P(\|dB'\|^2+\|\bQ B'\|^2)$. Therefore
\begin{align*}
\|B'\|^2+\|\dl g\|^2&\le C_P\|dB'\|^2+(C_P+2)\big(\|\bQ B'\|^2+\tfrac12\|\dl g\|^2\big)\\
&\le\big(C_P+2C_W^2(C_P+2)\big)\|dB'\|^2+(C_P+2)\|\bQ B'+\dl g\|^2 . \qedhere
\end{align*}
\end{proof}

\subsection{Coercivity}

\begin{lemma}[Gradient part in Ba\l aban's gauge]\label{lem:gradient-part}
Let $\lambda$ be a scalar field with mean zero and $R(-\Lap\lambda)=0$. Then $\|d\lambda\|\le C_B\|\dl Q_0\lambda\|$, with $C_B$ from Lemma~\ref{lem:biharmonic}.
\end{lemma}

\begin{proof}
$-\Lap\lambda$ lies in $\cW^\perp=\cV$ and has mean zero. By Proposition~\ref{prop:R}(c), $-\Lap\lambda=(-\Lap)^{-1}Q_0^*\omega$ with $\omega$ of mean zero. Since $\lambda$ has mean zero, $\lambda=(-\Lap)^{-2}Q_0^*\omega$, and Lemma~\ref{lem:biharmonic} applies.
\end{proof}

\begin{theorem}[Coercivity of the gauge-field operator]\label{thm:gauge-coercive}
There is $c_g>0$, depending only on $d$ and $a$, such that for every gauge field $A$ on $\Tf$
\[
\ip A{\bL_aA}=\|dA\|^2+\|Rd^*A\|^2+a\|\bQ A\|^2\ \ge\ c_g\big(\|A\|^2+\|\partial A\|^2\big).
\]
\end{theorem}

\begin{proof}
\emph{Step 1 (gauge fixing).} Let $\lambda_1\in\cN(Q_0)$ be given by Lemma~\ref{lem:gauge-fix} and $B=A+d\lambda_1$. Then $Rd^*B=0$, $dB=dA$, $\bQ B=\bQ A$, and $\|d\lambda_1\|\le\sqrt{d/2}\|Rd^*A\|$.

\emph{Step 2 (Hodge decomposition).} By Lemma~\ref{lem:hodge}, $B=d\lambda+B'$ with $\lambda=(-\Lap)^{-1}d^*B$ of mean zero and $d^*B'=0$. Then $R(-\Lap\lambda)=Rd^*B=0$. By Lemma~\ref{lem:gradient-part}, $\|d\lambda\|\le C_B\|\dl g\|$ with $g=Q_0\lambda$.

\emph{Step 3 (separation).} By gauge covariance, $\bQ B=\bQ B'+\dl g$, and $dB'=dB=dA$. Lemma~\ref{lem:separation} gives
\[
\|B'\|^2+\|\dl g\|^2\le C_\star\big(\|dA\|^2+\|\bQ A\|^2\big).
\]

\emph{Step 4.} Since $A=-d\lambda_1+d\lambda+B'$,
\[
\|A\|\le\sqrt{d/2}\,\|Rd^*A\|+(C_B+1)\sqrt{2C_\star}\,\big(\|dA\|^2+\|\bQ A\|^2\big)^{1/2}.
\]
Hence $\|A\|^2\le C'\big(\|dA\|^2+\|Rd^*A\|^2+\|\bQ A\|^2\big)$.

\emph{Step 5 (derivatives).} By Lemma~\ref{lem:hodge}, $\|\partial A\|^2=\|dA\|^2+\|d^*A\|^2$. Moreover $\|d^*A\|\le\|Rd^*A\|+\|P_\cV d^*A\|\le\|Rd^*A\|+C\|A\|$ by Proposition~\ref{prop:R}(e). Combining with Step~4 proves the theorem, with $c_g$ depending on $\min(1,a)$ and the constants above.
\end{proof}

\begin{corollary}\label{cor:X-positive}
On the space $\cX=\{A:\bQ A=0,\ Rd^*A=0\}$ one has $\|dA\|^2\ge c_g\|A\|^2$.
\end{corollary}

This completes the proof of Theorem~\ref{thm:A}.
